\documentclass[10pt,journal,compsoc,onecolumn]{IEEEtran}
\usepackage[utf8]{inputenc}
\usepackage{amsmath,amsfonts,amssymb}
\usepackage{amsthm}
\usepackage{graphicx}
\usepackage{cite}
\usepackage{textcomp}
\usepackage{tikz}
\usetikzlibrary{arrows,positioning,shapes.geometric}
\usepackage{booktabs}
\usepackage{url}

\newtheorem{theorem}{Theorem}
\newtheorem{lemma}[theorem]{Lemma}
\newtheorem{definition}[theorem]{Definition}
\newtheorem{corollary}[theorem]{Corollary}

\begin{document}

\title{System Design Masterclass Vol. 1: Low-Level Design (LLD) --- OOP Foundations, SOLID Principles, and Core GoF Design Patterns in Java}

\IEEEtitleabstractindextext{%
\begin{abstract}
A definitive, practice-oriented engineering masterclass on Low-Level Design (LLD) and Machine Coding in Java, designed specifically for a 30-day intensive mastery curriculum (Days 1--7). This volume bridges the gap between academic theory and real-world software craftsmanship by eliminating abstract mathematical bloat and focusing entirely on actionable principles: the four OOP pillars from a clean-code perspective, rigorous before-and-after refactoring across all five SOLID principles, and complete runnable Java reference implementations for the top seven Gang of Four (GoF) design patterns (Strategy, Observer, Factory, Decorator, Thread-Safe Singleton with Double-Checked Locking, Builder, and State). It concludes with a structured 90-minute machine coding interview execution framework and an actionable daily 2-hour study schedule.
\end{abstract}

\begin{IEEEkeywords}
binary trees, algorithms, data structures, computational complexity, tree traversal, recursive algorithms, formal analysis
\end{IEEEkeywords}
}

\maketitle
\IEEEdisplaynontitleabstractindextext
\IEEEpeerreviewmaketitle

\section{The 30-Day System Design Master Roadmap (2 Hours/Day)}
\label{sec:roadmap}

Mastering System Design in 30 days requires an aggressive, structured, and practical approach. Instead of reading theoretical treatises, this 30-day curriculum allocates exactly 2 hours each day (60 hours total), divided equally between **Low-Level Design (LLD / Machine Coding in Java)** and **High-Level Design (HLD / Distributed Systems Architecture)**.

\begin{table}[htbp]
\caption{The 30-Day System Design Intensive Curriculum (2 Hours / Day)}
\label{tab:30day_roadmap}
\centering
\begin{tabular}{cp{4.5cm}p{6.0cm}l}
\toprule
\textbf{Days} & \textbf{Module / Volume} & \textbf{Core Focus \& Key Topics} & \textbf{Daily Commitment} \\
\midrule
\textbf{Days 1--7} & \textbf{Vol. 1: LLD Core Foundations} & OOP Pillars, SOLID Principles, Top 7 GoF Patterns with runnable Java implementations. & 2 hrs / day (14 hrs) \\
\midrule
\textbf{Days 8--15} & \textbf{Vol. 2: LLD Machine Coding \& Concurrency} & Multi-threading (`ReentrantLock`, `ExecutorService`), Parking Lot, BookMyShow, Splitwise, LRU Cache. & 2 hrs / day (16 hrs) \\
\midrule
\textbf{Days 16--22} & \textbf{Vol. 3: HLD Distributed Primitives} & Load Balancers, Consistent Hashing, Caching (Redis), SQL vs NoSQL, Sharding, Kafka, CAP/PACELC. & 2 hrs / day (14 hrs) \\
\midrule
\textbf{Days 23--30} & \textbf{Vol. 4: HLD Enterprise Case Studies} & TinyURL, WhatsApp/Chat, Twitter Feed, Uber Geohash, Amazon Flash Sale, 45-min Interview Framework. & 2 hrs / day (16 hrs) \\
\bottomrule
\end{tabular}
\end{table}

This document constitutes **Volume 1: Low-Level Design (LLD) Core Foundations**, covering Days 1 through 7.

% -----------------------------------------------------------------------------
\section{The Four OOP Pillars: An Interview \& Clean Code Perspective}
\label{sec:oop_pillars}

Every candidate can recite textbook definitions of OOP. In top-tier machine coding rounds (Amazon, Flipkart, Uber, Swiggy), interviewers evaluate how you apply these pillars to produce maintainable, extensible Java code.

\subsection{Encapsulation: Invariant Protection, Not Just Getters/Setters}
Encapsulation is not merely wrapping fields in `private` and generating getters and setters. True encapsulation guarantees that an object maintains its internal invariants and business rules.

\begin{verbatim}
// Bad: Leaky encapsulation (anemic model)
public class BankAccount {
    public double balance; // Anyone can set negative balance directly
}

// Good: Strict encapsulation with validated state transitions
public class BankAccount {
    private double balance;

    public BankAccount(double initialDeposit) {
        if (initialDeposit < 0) throw new IllegalArgumentException("Initial deposit >= 0");
        this.balance = initialDeposit;
    }

    public synchronized void deposit(double amount) {
        if (amount <= 0) throw new IllegalArgumentException("Deposit amount must be > 0");
        this.balance += amount;
    }

    public synchronized void withdraw(double amount) {
        if (amount <= 0 || amount > this.balance) {
            throw new IllegalStateException("Insufficient funds or invalid withdrawal amount");
        }
        this.balance -= amount;
    }

    public double getBalance() { return balance; }
}
\end{verbatim}

\subsection{Abstraction: Programming to an Interface}
Abstraction hides implementation complexity behind stable contracts (`interface` or `abstract class`). In machine coding, you must always program to interfaces. This decouples the caller from the underlying provider, making services interchangeable.

\subsection{Inheritance vs. Composition: Why Composition Wins}
A classic junior mistake is creating deep inheritance hierarchies (`Vehicle -> LandVehicle -> WheeledVehicle -> Car -> SportsCar`). Deep hierarchies lead to fragile base class problems, code duplication, and tight coupling.

\textbf{The Golden Rule of LLD:} \emph{Favor object composition over class inheritance.}
\begin{itemize}
    \item \textbf{Inheritance (IS-A):} Use only when true behavioral substitution holds under all circumstances.
    \item \textbf{Composition (HAS-A):} Combine independent behaviors via injected components. If a car needs an engine, a `Car` should \emph{have} an `Engine`, not \emph{inherit} from `Engine`.
\end{itemize}

\subsection{Polymorphism: Eliminating Brittle If-Else Chains}
Polymorphism allows treating different underlying objects through a uniform interface. In clean LLD, polymorphism is your primary weapon for eliminating ugly `switch` and `if-else` chains.

% -----------------------------------------------------------------------------
\section{Core Java OOP Language Mechanics \& Interview Essentials}
\label{sec:java_oop_mechanics}

While the four pillars provide the architectural theory, Java interviews frequently probe deep into the language-specific OOP mechanics that govern how the JVM executes object-oriented code.

\subsection{1. Abstract Class vs. Interface in Modern Java (Java 8 to 21)}
\begin{itemize}
    \item \textbf{State vs. Behavior:} An \texttt{abstract class} can define instance variables (mutable state) and constructors. An \texttt{interface} cannot hold mutable instance state (all fields are implicitly \texttt{public static final}).
    \item \textbf{Multiple Inheritance:} A Java class can implement multiple interfaces, but can extend only one abstract class (avoiding the diamond problem of multiple state inheritance).
    \item \textbf{Modern Java Interfaces (Java 8+):} Interfaces can define \texttt{default} methods (providing backward-compatible method bodies) and \texttt{static} utility methods. Java 9 added \texttt{private} interface helper methods.
    \item \textbf{When to use which:} Use an \texttt{interface} to define a peripheral capability or behavioral contract (\texttt{Flyable}, \texttt{Serializable}, \texttt{Comparable}). Use an \texttt{abstract class} when creating a tightly coupled family of classes that share common base code and mutable state (e.g., \texttt{BasePizza}, \texttt{AbstractGameEngine}).
\end{itemize}

\subsection{2. The \texttt{equals()} and \texttt{hashCode()} Contract}
One of the most frequently failed Java interview questions is explaining the strict mathematical contract between \texttt{equals()} and \texttt{hashCode()}:
\begin{enumerate}
    \item If two objects are equal according to \texttt{equals(Object)}, their \texttt{hashCode()} \textbf{must} produce the identical integer.
    \item If two objects have the same \texttt{hashCode()}, they are \textbf{not} necessarily equal (hash collision).
    \item If you override \texttt{equals()}, you \textbf{must} override \texttt{hashCode()}.
\end{enumerate}

\textbf{What breaks if you fail this rule?}
If you insert a custom object into a \texttt{HashSet} or use it as a key in a \texttt{HashMap}, Java uses \texttt{hashCode()} to compute the bucket array index. If you override \texttt{equals()} but forget \texttt{hashCode()}, two logically identical objects land in completely different buckets! As a result, \texttt{map.get(key)} returns \texttt{null}, and \texttt{set.contains(key)} fails silently, causing catastrophic state corruption in production!

\subsection{3. Compile-Time vs. Runtime Polymorphism}
\begin{itemize}
    \item \textbf{Compile-Time (Static / Method Overloading):} Methods share the same name in the same class but have different parameter signatures (type, count, or order). Resolved at compile-time by the compiler. Return type differences alone are not sufficient.
    \item \textbf{Runtime (Dynamic / Method Overriding):} A subclass provides a specific implementation of a method declared in its parent class. Resolved at runtime by the JVM via \textbf{Dynamic Method Dispatch} using the virtual method table (vtable).
\end{itemize}

\subsection{4. The \texttt{final} Keyword \& The Immutable Class Pattern}
In concurrent Java applications, immutability is the ultimate defense against race conditions. A class is truly immutable if its state cannot be modified after construction:
\begin{enumerate}
    \item Declare the class \texttt{final} (prevents subclassing).
    \item Make all fields \texttt{private} and \texttt{final}.
    \item Do not provide setter methods.
    \item \textbf{Defensive Copying:} If fields reference mutable objects (e.g., \texttt{Date}, \texttt{List}, \texttt{Map}), perform a deep copy in both the constructor and getters:
\end{enumerate}

\begin{verbatim}
public final class ImmutablePosition {
    private final int x;
    private final int y;
    private final List<String> tags;

    public ImmutablePosition(int x, int y, List<String> tags) {
        this.x = x;
        this.y = y;
        this.tags = new ArrayList<>(tags); // Defensive copy in constructor!
    }

    public List<String> getTags() {
        return Collections.unmodifiableList(tags); // Defensive copy in getter!
    }
}
\end{verbatim}

\subsection{5. Java Access Modifiers Matrix}
\begin{table}[htbp]
\caption{Java Access Modifiers Scope Matrix}
\label{tab:access_modifiers}
\centering
\begin{tabular}{lcccc}
\toprule
\textbf{Modifier} & \textbf{Same Class} & \textbf{Same Package} & \textbf{Subclass (Different Pkg)} & \textbf{World (Anywhere)} \\
\midrule
\texttt{public} & Yes & Yes & Yes & Yes \\
\texttt{protected} & Yes & Yes & Yes & No \\
\texttt{default} (Package-Private) & Yes & Yes & No & No \\
\texttt{private} & Yes & No & No & No \\
\bottomrule
\end{tabular}
\end{table}

\subsection{6. Constructor Chaining: \texttt{this()} and \texttt{super()}}
\begin{itemize}
    \item \texttt{this(...)} invokes an overloaded constructor within the same class.
    \item \texttt{super(...)} invokes the direct superclass constructor.
    \item \textbf{Rule:} Either \texttt{this()} or \texttt{super()} must be the \textbf{first line} in a constructor body, ensuring the object hierarchy initializes from root \texttt{Object} down to the leaf instance.
\end{itemize}

\subsection{7. Comparable vs. Comparator (Sorting \& Ordering)}
\begin{itemize}
    \item \texttt{Comparable<T>} \textbf{(Internal / Natural Ordering):} Implemented by the domain class itself via \texttt{compareTo(T o)}. Defines the single canonical sort order.
    \item \texttt{Comparator<T>} \textbf{(External / Custom Strategy):} Implemented as a standalone class or lambda \texttt{(a, b) -> a.val - b.val}. Allows defining multiple independent sorting strategies (e.g., sort by price, sort by name, sort by date) without altering the entity.
\end{itemize}

\begin{verbatim}
// Natural ordering: by rating
public class Player implements Comparable<Player> {
    private final String name;
    private final int elo;

    public Player(String name, int elo) { this.name = name; this.elo = elo; }
    public int getElo() { return elo; }
    public String getName() { return name; }

    @Override
    public int compareTo(Player other) { return Integer.compare(this.elo, other.elo); }
}

// Custom strategy: sort alphabetically by name
Comparator<Player> byName = Comparator.comparing(Player::getName);
\end{verbatim}

\subsection{8. Fail-Fast vs. Fail-Safe Iterators}
\begin{itemize}
    \item \textbf{Fail-Fast (\texttt{ArrayList}, \texttt{HashMap}, \texttt{HashSet}):} Maintain an internal \texttt{modCount} counter. If the collection is structurally modified (item added or deleted) while an iterator is traversing it, the iterator immediately throws \texttt{ConcurrentModificationException}. \textbf{Rule:} To remove during iteration, always use \texttt{iterator.remove()}, never \texttt{list.remove()}!
    \item \textbf{Fail-Safe (\texttt{ConcurrentHashMap}, \texttt{CopyOnWriteArrayList}):} Operates on a snapshot or clone of the underlying collection. Structural mutations do not throw exceptions, guaranteeing safe multi-threaded iteration at the cost of slight memory overhead.
\end{itemize}

% -----------------------------------------------------------------------------
\section{The SOLID Principles Masterclass (With Real Java Refactoring)}
\label{sec:solid}

The SOLID principles are the ultimate design guidelines for writing decoupled, testable, and enterprise-grade software.

\subsection{S --- Single Responsibility Principle (SRP)}
\textbf{Rule:} \emph{A class should have one, and only one, reason to change.}

\textbf{Anti-Pattern (The God Object):}
\begin{verbatim}
// Bad: Invoice does business math, database persistence, and PDF generation!
public class Invoice {
    private List<Item> items;

    public double calculateTotal() { /* business math */ return 0.0; }
    public void saveToDatabase() { /* JDBC/SQL persistence logic */ }
    public void printPdfReceipt() { /* PDF rendering logic */ }
}
\end{verbatim}

\textbf{Refactored Solution (Strict SRP):}
\begin{verbatim}
public class Invoice {
    private final List<Item> items;
    public Invoice(List<Item> items) { this.items = items; }
    public double calculateTotal() {
        return items.stream().mapToDouble(Item::getPrice).sum();
    }
}

public class InvoiceRepository {
    public void save(Invoice invoice) {
        // Responsible ONLY for database persistence
    }
}

public class InvoicePrinter {
    public void print(Invoice invoice) {
        // Responsible ONLY for presentation / printing
    }
}
\end{verbatim}

\subsection{O --- Open/Closed Principle (OCP)}
\textbf{Rule:} \emph{Software entities should be open for extension, but closed for modification.}

Whenever you add a new business requirement (e.g., adding UPI or Crypto payments), you should **never modify existing, tested code**. Instead, add a new class implementing a shared interface.

\textbf{Refactored Java Solution (Strategy Pattern):}
\begin{verbatim}
public interface PaymentStrategy {
    void pay(double amount);
}

public class CreditCardPayment implements PaymentStrategy {
    public void pay(double amount) { System.out.println("Paid " + amount + " via Credit Card"); }
}

public class UpiPayment implements PaymentStrategy {
    public void pay(double amount) { System.out.println("Paid " + amount + " via UPI"); }
}

// Open for new payment methods without changing PaymentProcessor
public class PaymentProcessor {
    public void processPayment(PaymentStrategy strategy, double amount) {
        strategy.pay(amount);
    }
}
\end{verbatim}

\subsection{L --- Liskov Substitution Principle (LSP)}
\textbf{Rule:} \emph{Subtypes must be substitutable for their base types without altering program correctness.}

If class $S$ is a subclass of class $T$, any code expecting an instance of $T$ must work seamlessly with an instance of $S$ without throwing unexpected runtime exceptions or violating invariants.

\textbf{The Classic Anti-Pattern: Flying Bird vs Penguin:}
\begin{verbatim}
// Bad: Penguin violates the contract of Bird!
public class Bird {
    public void fly() { System.out.println("Flying..."); }
}

public class Penguin extends Bird {
    @Override
    public void fly() {
        throw new UnsupportedOperationException("Penguins cannot fly!"); // Violates LSP!
    }
}
\end{verbatim}

\textbf{Refactored Solution:}
\begin{verbatim}
public interface Bird {
    void eat();
}

public interface FlyingBird extends Bird {
    void fly();
}

public class Sparrow implements FlyingBird {
    public void eat() { System.out.println("Eating seeds"); }
    public void fly() { System.out.println("Flying high"); }
}

public class Penguin implements Bird {
    public void eat() { System.out.println("Eating fish"); }
    // No fly() method -> LSP perfectly preserved!
}
\end{verbatim}

\subsection{I --- Interface Segregation Principle (ISP)}
\textbf{Rule:} \emph{Clients should not be forced to depend upon interfaces that they do not use.}

Break "fat" interfaces into smaller, role-specific interfaces.

\textbf{Refactored Solution:}
\begin{verbatim}
// Bad: Fat interface forces unnecessary implementations
public interface Worker {
    void work();
    void eat();
    void sleep();
}

// Good: Segregated role-based interfaces
public interface Workable { void work(); }
public interface Feedable { void eat(); }

public class HumanWorker implements Workable, Feedable {
    public void work() { System.out.println("Human working"); }
    public void eat() { System.out.println("Human eating lunch"); }
}

public class RobotWorker implements Workable {
    public void work() { System.out.println("Robot assembling parts 24/7"); }
    // No need to implement eat()!
}
\end{verbatim}

\subsection{D --- Dependency Inversion Principle (DIP)}
\textbf{Rule:} \emph{High-level modules should not depend on low-level modules; both should depend on abstractions.}

\textbf{Refactored Solution (Dependency Injection):}
\begin{verbatim}
public interface MessageSender {
    void send(String recipient, String message);
}

public class EmailSender implements MessageSender {
    public void send(String recipient, String message) {
        System.out.println("Sending Email to " + recipient + ": " + message);
    }
}

public class SmsSender implements MessageSender {
    public void send(String recipient, String message) {
        System.out.println("Sending SMS to " + recipient + ": " + message);
    }
}

// High-level service depends on abstraction (MessageSender), not concrete EmailSender
public class NotificationService {
    private final MessageSender sender;

    // Injected via constructor
    public NotificationService(MessageSender sender) {
        this.sender = sender;
    }

    public void notifyUser(String user, String alert) {
        sender.send(user, alert);
    }
}
\end{verbatim}

% -----------------------------------------------------------------------------
\section{Top 7 GoF Design Patterns for Machine Coding Interviews}
\label{sec:design_patterns}

Of the 23 Gang of Four (GoF) patterns, these 7 appear in 95\% of all machine coding interview rounds.

\subsection{1. Strategy Pattern (Behavioral)}
\textbf{Intent:} Defines a family of algorithms, encapsulates each one, and makes them interchangeable at runtime.
\textbf{Common Interview Scenarios:} Discount engines, payment processing, route planning (fastest, cheapest, scenic), sorting strategies.

\subsection{2. Observer Pattern (Behavioral)}
\textbf{Intent:} Defines a one-to-many dependency between objects so that when one object changes state, all its dependents are notified automatically.
\textbf{Common Interview Scenarios:} Stock market ticker, YouTube subscriber alerts, Cricbuzz live cricket score updates, in-memory event buses.

\begin{verbatim}
import java.util.*;

public interface Observer {
    void update(String score);
}

public interface Subject {
    void register(Observer o);
    void unregister(Observer o);
    void notifyObservers();
}

public class CricketScoreBoard implements Subject {
    private final List<Observer> observers = new ArrayList<>();
    private String currentScore;

    public void setScore(String score) {
        this.currentScore = score;
        notifyObservers();
    }

    public void register(Observer o) { observers.add(o); }
    public void unregister(Observer o) { observers.remove(o); }
    public void notifyObservers() {
        for (Observer o : observers) { o.update(currentScore); }
    }
}
\end{verbatim}

\subsection{3. Factory \& Factory Method Pattern (Creational)}
\textbf{Intent:} Defines an interface for creating an object, but lets subclasses decide which class to instantiate. Decouples object instantiation from business logic.
\textbf{Common Interview Scenarios:} Vehicle allocation in a Parking Lot, database driver connections, UI widget renderers.

\begin{verbatim}
public interface Vehicle { void park(); }
public class Car implements Vehicle { public void park() { System.out.println("Car parked"); } }
public class Bike implements Vehicle { public void park() { System.out.println("Bike parked"); } }

public class VehicleFactory {
    public static Vehicle createVehicle(String type) {
        if ("CAR".equalsIgnoreCase(type)) return new Car();
        if ("BIKE".equalsIgnoreCase(type)) return new Bike();
        throw new IllegalArgumentException("Unknown vehicle type: " + type);
    }
}
\end{verbatim}

\subsection{4. Decorator Pattern (Structural)}
\textbf{Intent:} Attaches additional responsibilities to an object dynamically at runtime without modifying the underlying class (alternative to subclassing).
\textbf{Common Interview Scenarios:} Pizza toppings (Base Pizza + Extra Cheese + Mushroom), Coffee maker, Java I/O streams (`BufferedReader(new FileReader())`).

\begin{verbatim}
public interface BasePizza { double cost(); }
public class Margherita implements BasePizza { public double cost() { return 150.0; } }

public abstract class PizzaToppingDecorator implements BasePizza {
    protected final BasePizza pizza;
    public PizzaToppingDecorator(BasePizza pizza) { this.pizza = pizza; }
}

public class ExtraCheese extends PizzaToppingDecorator {
    public ExtraCheese(BasePizza pizza) { super(pizza); }
    public double cost() { return pizza.cost() + 40.0; }
}

public class Jalapeno extends PizzaToppingDecorator {
    public Jalapeno(BasePizza pizza) { super(pizza); }
    public double cost() { return pizza.cost() + 30.0; }
}

// Client: BasePizza myPizza = new Jalapeno(new ExtraCheese(new Margherita()));
// Total Cost = 150 + 40 + 30 = 220.0
\end{verbatim}

\subsection{5. Thread-Safe Singleton Pattern (Creational)}
\textbf{Intent:} Ensures a class has only one instance and provides a global access point to it.
\textbf{Interview Requirement:} You must write the **Double-Checked Locking** version using `volatile` to prevent CPU instruction reordering.

\begin{verbatim}
public class DatabaseConnectionPool {
    // volatile is MANDATORY to prevent partially constructed object leakage
    private static volatile DatabaseConnectionPool instance;

    private DatabaseConnectionPool() {
        // Prevent reflection instantiation attack
        if (instance != null) throw new RuntimeException("Use getInstance() method");
    }

    public static DatabaseConnectionPool getInstance() {
        if (instance == null) { // First check (no lock)
            synchronized (DatabaseConnectionPool.class) {
                if (instance == null) { // Second check (with lock)
                    instance = new DatabaseConnectionPool();
                }
            }
        }
        return instance;
    }
}
\end{verbatim}

\subsection{6. Builder Pattern (Creational)}
\textbf{Intent:} Separates the construction of a complex object from its representation, allowing the same construction process to create different representations.
\textbf{Interview Requirement:} Solves the "Telescoping Constructor" anti-pattern when an object has 5+ optional configuration parameters.

\begin{verbatim}
public class UserProfile {
    private final String userId;     // Required
    private final String email;      // Required
    private final String phone;      // Optional
    private final int age;           // Optional

    private UserProfile(Builder b) {
        this.userId = b.userId;
        this.email = b.email;
        this.phone = b.phone;
        this.age = b.age;
    }

    public static class Builder {
        private final String userId;
        private final String email;
        private String phone = "";
        private int age = 0;

        public Builder(String userId, String email) {
            this.userId = userId;
            this.email = email;
        }

        public Builder phone(String phone) { this.phone = phone; return this; }
        public Builder age(int age) { this.age = age; return this; }
        public UserProfile build() { return new UserProfile(this); }
    }
}
// Usage: UserProfile user = new UserProfile.Builder("101", "pavan@work.com").age(28).build();
\end{verbatim}

\subsection{7. State Pattern (Behavioral)}
\textbf{Intent:} Allows an object to alter its behavior when its internal state changes. The object will appear to change its class.
\textbf{Common Interview Scenarios:} Vending Machine, Order Delivery Lifecycle (Placed $\to$ Paid $\to$ Shipped $\to$ Delivered $\to$ Returned).

\subsection{8. Command Pattern (Behavioral) --- Undo/Redo \& Decoupled Actions}
\textbf{Intent:} Encapsulates a request as an object, thereby letting you parameterize clients with different requests, queue or log requests, and support \textbf{reversible operations (Undo/Redo)}.
\textbf{Critical Interview Use Cases:} Chess Engine move history (Undo/Redo), transactional database rollbacks, graphical editor operations.

\begin{verbatim}
import java.util.*;

public interface Command {
    void execute();
    void undo();
}

public class MovePieceCommand implements Command {
    private final Board board;
    private final Position from;
    private final Position to;
    private Piece capturedPiece;

    public MovePieceCommand(Board board, Position from, Position to) {
        this.board = board; this.from = from; this.to = to;
    }

    @Override
    public void execute() {
        this.capturedPiece = board.getPiece(to);
        board.move(from, to);
    }

    @Override
    public void undo() {
        board.move(to, from); // Reverse move
        if (capturedPiece != null) board.placePiece(to, capturedPiece); // Restore captured piece
    }
}

public class CommandHistory {
    private final Deque<Command> undoStack = new ArrayDeque<>();
    private final Deque<Command> redoStack = new ArrayDeque<>();

    public void executeCommand(Command cmd) {
        cmd.execute();
        undoStack.push(cmd);
        redoStack.clear(); // Clear redo stack on new action
    }

    public void undo() {
        if (!undoStack.isEmpty()) {
            Command cmd = undoStack.pop();
            cmd.undo();
            redoStack.push(cmd);
        }
    }
}
\end{verbatim}

% -----------------------------------------------------------------------------
\section{The 90-Minute Machine Coding Interview Execution Framework}
\label{sec:framework}

In top tech interviews, you are given 90 minutes to design and implement a fully functioning, compilable low-level system in your IDE. Follow this exact timed blueprint:

\begin{table}[htbp]
\caption{The 90-Minute Machine Coding Time Management Blueprint}
\label{tab:lld_blueprint}
\centering
\begin{tabular}{ccl}
\toprule
\textbf{Time Interval} & \textbf{Phase} & \textbf{Action Items \& Key Deliverables} \\
\midrule
0 -- 10 mins & Requirement Discovery & Clarify functional scope, identify boundary constraints, list core actors. \\
10 -- 25 mins & Class Diagram \& Schema & Design domain entities (`models`), interfaces, and service contracts. \\
25 -- 65 mins & Clean Implementation & Code models, strategies, factories, and service business logic in Java. \\
65 -- 75 mins & Concurrency \& Validation & Add thread-safety (`synchronized`/`ReentrantLock`) and input validations. \\
75 -- 90 mins & Driver Code \& Testing & Write a rich `Main.java` demonstrating all use cases and edge cases. \\
\bottomrule
\end{tabular}
\end{table}

\subsection{Production Error Handling \& Custom Exception Hierarchy}
Interviewers severely penalize candidates who return \texttt{null} or boolean flags (\texttt{return false;}) when business rules fail. Production machine coding requires a structured unchecked exception hierarchy:

\begin{verbatim}
// Base domain exception (Unchecked RuntimeException)
public abstract class DomainException extends RuntimeException {
    public DomainException(String message) { super(message); }
}

public class ResourceNotFoundException extends DomainException {
    public ResourceNotFoundException(String resource, String id) {
        super(resource + " not found with ID: " + id);
    }
}

public class InvalidOperationException extends DomainException {
    public InvalidOperationException(String message) { super(message); }
}

// In service logic:
// if (spot == null) throw new ResourceNotFoundException("ParkingSpot", spotId);
// if (!spot.isFree()) throw new InvalidOperationException("Spot already occupied!");
\end{verbatim}

% -----------------------------------------------------------------------------
\section{Daily 2-Hour Action Plan for Days 1 to 7}
\label{sec:action_plan}

\begin{itemize}
    \item \textbf{Day 1 (2 hrs):} Read Section 2 (OOP Pillars). Implement the Encapsulated `BankAccount` and composition examples in your IDE.
    \item \textbf{Day 2 (2 hrs):} Master Single Responsibility (SRP) and Open/Closed (OCP). Refactor an existing `switch-case` discount calculator to the Strategy pattern.
    \item \textbf{Day 3 (2 hrs):} Master Liskov Substitution (LSP), Interface Segregation (ISP), and Dependency Inversion (DIP). Implement Dependency Injection for a notification engine.
    \item \textbf{Day 4 (2 hrs):} Code the **Strategy Pattern** and **Observer Pattern** from scratch. Build an in-memory Stock Price Notification simulator.
    \item \textbf{Day 5 (2 hrs):} Code the **Factory Method** and **Decorator Pattern**. Build a customizable Pizza billing calculator with dynamic toppings.
    \item \textbf{Day 6 (2 hrs):} Implement the **Thread-Safe Double-Checked Locking Singleton** and the **Builder Pattern**. Test against multi-threaded access.
    \item \textbf{Day 7 (2 hrs):} Implement a mini-project (e.g., A Vending Machine or State-driven Order Pipeline) combining Factory, Strategy, and State patterns.
\end{itemize}


\end{document}
